\chapter{Scientific Background\label{cha:background}}
% Drafted per POINT-SEQUENCE.md Ch.2 (points 2.1–2.20). Point IDs marked per paragraph.
% Style: thesis-prose (one idea per sentence, present tense, no term before its definition).

This chapter introduces the systems and terms used throughout the rest of the thesis.
It first describes NebulaStream and its execution model.
It then explains how NebulaStream computes aggregates over time windows.
It continues with RocksDB, the storage engine used in this thesis.
It closes with the principle behind moving data to storage and with the metrics used in the evaluation.

\section{Stream Processing and NebulaStream\label{sec:bg-nes}}

A data stream is a continuous real-time transmission of data.
Each data item is called a tuple.
A tuple is one record, for example, one sensor reading together with the time it was taken.
A stream has no defined end.
A \ac{DSPS} computes results in real time as data arrives from such streams.
It does not first store all the data and query it later for evaluation.

NebulaStream is a data management system for the \ac{IoT}~\cite{zeuch2020,zeuch2022}.
Researchers at \ac{DIMA} develop it together with \ac{BIFOLD}.
Its execution engine processes data streams; this thesis builds on that engine.
The \ac{IoT} consists of many small networked devices, such as sensors spread across a city. 
This part of the network is called the edge.
NebulaStream combines the advantages of edge processing near the source, such as low latency and reduced bandwidth, with those of cloud processing in data centers, such as powerful computation and storage~\cite{zeuch2020,zeuch2022}.
% [2.1]
A user describes a computation as a query.
NebulaStream compiles each query into pipelines~\cite{zeuch2020,zeuch2022}.
A pipeline is a chain of operators.
An operator is a single processing step, for example, a filter or a sum.

% [2.2]
Execution is push-based.
Sources push tuples through the pipelines as data arrives.
No part of the system polls for new data~\cite{zeuch2020,zeuch2022}.

% [2.3]
A fixed set of worker threads executes the pipeline work.
The work is split into tasks.
The worker threads take tasks from a shared queue.

% [2.4]
Tuples are transmitted and processed in buffers.
A buffer is a block of memory with a fixed size.

% [2.5]
All buffers come from one global buffer pool.
The pool holds a fixed number of buffers.
The engine allocates the pool once at process start.
This pool is the central shared resource in this thesis.

% [2.6]
A worker thread that needs a buffer requests one from the pool.
If the pool is empty, the thread waits up to one second.
If no buffer becomes free in that time, the engine aborts the query.

Queries that group data over time must remember intermediate results between tuples.
The next section explains how NebulaStream organizes this memory.

\section{Windowed Aggregation\label{sec:bg-windows}}

An aggregate summarizes many tuples into a single value.
Examples are a count, a sum, or an average.
A stream has no end, so an aggregation over the whole stream never finishes.
Windows solve this problem.

% [2.7]
A tumbling window splits the stream into time intervals of fixed length.
The intervals do not overlap.
The operator computes one aggregate per key and per window.
A key identifies one group in the data, for example one sensor.

% [2.8]
As tuples arrive, a window accumulates aggregation values.
These values are called state.
NebulaStream stores window state using stream slicing~\cite{TraubGCBKRM20}.
A slice is a time interval that holds the aggregate of the tuples that fall in it.
A new slice starts at every window boundary, and slices never overlap.
A window's result is assembled by combining the slices that fall inside it.
When windows overlap, each window is built from several slices, and a slice can serve multiple windows.
This keeps every tuple aggregated only once.
Tumbling windows do not overlap, so each window is built from exactly one slice.
That single slice holds the window's entire state.

% [2.9]
Inside a slice, the state is stored in a hash map.
A hash map is a data structure that finds the entry for a key in constant time.
NebulaStream uses a hash map implementation called ChainedHashMap.
Its entries live on pages.
A page is a fixed-size block of memory that holds hash-map entries.
A separate table, the bucket directory, points to the entries.
Each entry consists of a 24-byte header and the aggregate values themselves, the payload.

% [2.10 — binding vocabulary]
The pages of a hash map come from one of two sources.
With \emph{bounded allocation}, pages are taken from the fixed global buffer pool and are returned to it on release.
With \emph{unbounded allocation}, pages come from the process's general-purpose memory, the heap.
Heap memory is limited only by what the operating system grants the process.
Which of the two allocation types is chosen depends on the configured page size relative to the engine's configured buffer size.
Chapter~\ref{cha:evaluation} measures the differences between these two types of allocation.
Figure~\ref{fig:nes-components} summarizes the components introduced so far and where window state lives among them.

% Background view: the unmodified NebulaStream components that hold window state.
% Deliberately contains NO thesis extensions (no spiller, no RocksDB, no thresholds):
% those appear in fig:system-overview (Section 4.1), which extends this picture.
% Same visual language as the Chapter-4 figures; everything here is existing engine
% code, so everything is grey.

\begin{figure}[ht]
\centering
\resizebox{\textwidth}{!}{%
\begin{tikzpicture}[
  font=\sffamily\small,
  every node/.style={align=center},
  box/.style     ={rectangle, draw, thick, minimum height=1.5cm, inner sep=5pt, rounded corners=1pt, minimum width=3.1cm},
  existing/.style={box, fill=black!12},
  slice/.style   ={rectangle, draw, thick, minimum width=2.9cm, minimum height=1.5cm, inner sep=4pt, rounded corners=1pt, fill=black!7},
  flow/.style    ={-{Latex[length=2.6mm]}, very thick},
  watch/.style   ={-{Latex[length=2.4mm]}, thick, dashed, black!60},
  lbl/.style     ={fill=white, draw=black!30, line width=0.2pt, rounded corners=1pt,
                   inner sep=2.5pt, font=\sffamily\footnotesize\itshape, text=black!75},
]

% ---- row A: the processing path ----
\node[existing, minimum width=2.9cm] (src) at (0,0)
  {Sensor stream\\\footnotesize\itshape unbounded tuples};
\node[existing] (workers) at (4.6,0)
  {Worker threads\\\footnotesize\itshape take tasks from\\\footnotesize\itshape a shared queue};
\node[slice] (sliceA) at (9.4,0)
  {Slice (one window)\\\footnotesize\itshape hash maps\\\footnotesize\itshape on pages};
\node[slice] (sliceB) at (12.8,0)
  {Older slices\\\footnotesize\itshape complete, waiting\\\footnotesize\itshape for emission};
\node[existing, minimum width=3.0cm] (emit) at (17.6,0)
  {Emission\\\footnotesize\itshape results downstream;\\\footnotesize\itshape slices freed after};

% ---- slice store group box ----
\begin{pgfonlayer}{background}
  \node[draw, thick, rounded corners, fit=(sliceA)(sliceB), inner sep=4.5mm, fill=black!4] (store) {};
\end{pgfonlayer}
\node[anchor=south, font=\sffamily\footnotesize\bfseries, text=black!65]
  at (store.north) {Slice store \textnormal{\footnotesize\itshape (one slice per tumbling window)}};

% ---- row B: memory suppliers + event time ----
\node[existing, minimum width=4.6cm] (pool) at (4.6,-4.0)
  {Global buffer pool\\\footnotesize\itshape fixed number of buffers,\\\footnotesize\itshape allocated once at start};
\node[existing, minimum width=3.0cm] (heap) at (0,-4.0)
  {Process heap\\\footnotesize\itshape grows on demand};
\node[existing, minimum width=3.4cm] (wm) at (12.8,-4.0)
  {Watermark\\\footnotesize\itshape event time up to which\\\footnotesize\itshape all tuples have arrived};

% ---- edges ----
\draw[flow] (src) -- node[lbl]{tuples\\in buffers} (workers);
\draw[flow] (workers) -- node[lbl, pos=0.38]{writes\\(build)} (sliceA);
\draw[flow] (sliceA) -- node[lbl, below=3.5mm]{time passes} (sliceB);
\draw[flow] (sliceB.east) -- node[lbl, pos=0.5]{result per key\\and window} (emit.west);

\draw[flow] (pool.north) -- node[lbl, pos=0.7]{buffers} (workers.south);
\draw[flow] (pool.north east) -- node[lbl, pos=0.45]{pages (bounded\\allocation)} (sliceA.south west);
\draw[flow] (heap.north) -- ++(0,0.7) -| node[lbl, pos=0.08]{pages + directories\\(unbounded allocation)} (sliceA.south);
\draw[watch] (wm.north) -- node[lbl, pos=0.5]{marks windows\\complete} (sliceB.south);

\end{tikzpicture}%
}
\caption[NebulaStream components that hold window state]{%
The NebulaStream components that hold window state, before any change by this
thesis. Worker threads take tasks from a shared queue and write arriving tuples
into the slice of the current tumbling window. A slice stores its state in hash
maps, and the hash maps live on fixed-size pages. Pages come from the global
buffer pool when one page fits into one buffer (bounded allocation), and from
the process heap otherwise (unbounded allocation). Once the watermark passes a
window's end, the window is complete; its result is emitted downstream and its
slices are freed afterwards. All state lives in main memory. Chapter~\ref{cha:solution}
extends exactly this picture with a spillable slice store
(Figure~\ref{fig:system-overview}).}
\label{fig:nes-components}
\end{figure}


Worker threads take tasks from a shared queue and write arriving tuples
into the slice of the current tumbling window. A slice stores its state in hash
maps, and the hash maps live on fixed-size pages. Pages come from the global
buffer pool when one page fits into one buffer (bounded allocation), and from
the process heap otherwise (unbounded allocation). Once the watermark passes a
window's end, the window is complete; its result is emitted downstream, and its
slices are freed afterward. All state lives in main memory.
% [2.11]
Tuples can arrive out of order.
The watermark is the timestamp up to which all tuples are assumed to have arrived.
Windows are complete once the watermark has passed.

% [2.12]
A window passes through three phases.
While it fills (\texttt{WINDOW\_FILLING}), it still receives tuples.
Once the watermark passes the window's end, NebulaStream emits the result downstream and marks the window \texttt{EMITTED\_TO\_PROBE}; the result is now fixed, but the window's slices still occupy memory.
Garbage collection then deletes those slices and frees the memory after the result has been consumed downstream.

Window state grows with the number of keys and windows.
On a small device, it can grow beyond the available memory.
The system then needs to offload state to storage to prevent an \ac{OOM}-failure.
The next section introduces the storage system that provides this feature.

\section{RocksDB and the Log-Structured Merge Tree\label{sec:bg-rocksdb}}

A \ac{KV} store is a data structure that stores a value under a unique key.
It can quickly write, read, and delete the value for a given key.

% [2.13]
RocksDB is a persistent \ac{KV} store~\cite{borthakur2013}.
Persistent means that the data survives on storage.
RocksDB originates at Facebook and runs embedded inside another program.
Its design follows the \ac{LSM}~\cite{oneil1996}.

% [2.14]
Writes first go to a structure in memory, the memtable.
When the memtable is full, RocksDB writes its contents to a sorted file on storage and begins a new memtable.
These files accumulate, so a read may have to consult the memtable and several files to find the newest value for a key.

% [2.15]
A background process called compaction keeps this read cost bounded.
It merges sorted files into fewer, larger ones and discards values that a later write has overwritten or deleted.
Compaction therefore rewrites data that the application already wrote once.
The ratio of bytes written to storage over bytes written by the application is the write amplification~\cite{jin2022,wei2025}.

% [2.15b] — background-only survey of the parameters RocksDB exposes.
RocksDB exposes several parameters that control this behavior.
The write buffer size sets how large a memtable grows before RocksDB flushes it to storage.
The block cache size determines how much memory RocksDB caches data read from storage.
The compression algorithm determines how RocksDB encodes data in its files; stronger compression produces smaller files but requires more CPU work.
The compaction style determines how RocksDB merges files; the styles differ in read performance and write amplification.
A column family is a named partition of keys, with its own memtable and files, so that different data can use different settings.
These parameters together determine how much memory RocksDB uses~\cite{rocksdb-memory} and how much it writes~\cite{jin2022,wei2025}.

Prior studies show that RocksDB's parameters affect its write behavior and that tuning reduces, but does not eliminate, write amplification~\cite{jin2022,wei2025}.
Under a tight memory budget, these parameters must therefore be set deliberately; Section~\ref{sec:rocksdb-config} fixes the three that these studies identify, the compression, the write buffer size, and the block cache size, by measurement.

Storing state on storage is slower than keeping it in memory.
The next section explains when moving data to storage is still the right choice.

\section{The External-Memory Principle\label{sec:bg-extmem}}

% [2.16]
Algorithms whose data exceeds main memory are called external-memory algorithms~\cite{graefe1993}.
Their guiding principle is simple.
Data that is used often is hot.
Data that is used rarely is cold.
Keep hot data in memory and move cold data to storage.
The slower storage then holds only data that is rarely needed.

On a large server, memory for the hot data is usually available.
On an edge device, memory is scarce.
The machinery that moves data to and from storage also needs memory.
With a small fixed buffer pool, the hot data therefore includes the input and output path itself.
Chapter~\ref{cha:evaluation} returns to this point.

The evaluation measures memory and completion with a small set of terms.
The last section of this chapter defines them.

\section{Metrics and Terms\label{sec:bg-metrics}}

% [2.17]
Linux can place a group of processes under a hard memory limit.
The mechanism is called a \ac{cgroup}.
When the group reaches its limit, the kernel does not kill it at once.
It first reclaims memory: it drops or writes back the pages it is allowed to take away, to make room.
Only when this page reclaim cannot free enough does the kernel's OOM killer terminate the process.
This thesis writes OOM for out of memory.
In the evaluation, this cgroup is the one Docker creates for the container, so its hard limit is set by a Docker flag (Section~\ref{sec:reproducibility}); the bare term ``cgroup'' below always refers to that Docker-imposed control group.

% [2.17b] — how memory is measured, and why the distinction is load-bearing.
The kernel accounts a cgroup's memory in two kinds.
Anonymous memory holds the working data of the processes.
With swap disabled, as it is throughout this evaluation (Section~\ref{sec:reproducibility}), the kernel has nowhere to move an anonymous page, so it cannot reclaim it; anonymous memory is therefore the quantity that drives the OOM killer.
File-backed memory is the page cache for file reads and writes; the kernel reclaims it under pressure, so it never kills a run.
This distinction affects the measurement of a spilling backend in two ways.

First, a backend that spills to storage inflates the page cache as a by-product of working correctly, because every spilled and restored byte passes through it.
The kernel keeps these cached pages until something else needs the memory, because a clean cached page requires no work to keep and may save a future device read.
Under a hard limit, a spilling run's total charged memory therefore plateaus at the cap, held there by continuous page reclaim.
The page cache grows until anonymous memory and cache together fill the cap, and every further allocation is served by reclaiming cache.
Total consumption at the cap is thus the normal operating state of a spilling run, rather than a sign of imminent failure; the run fails only when the anonymous set has grown so large that the remaining cache cannot provide enough.

Second, a naive total, such as the resident set size or the cgroup's current consumption, counts the backend's own spill traffic and overstates its footprint in proportion to how much it spills.
Every memory footprint reported in this thesis is anonymous memory, read from the cgroup's memory statistics; the page cache is recorded separately and never counted as pressure.

% [2.18]
State pressure is measured as the ratio $\rho$.
$\rho$ is the steady-state anonymous-memory footprint of a workload divided by the memory cap.
A workload with $\rho > 1$ does not fit into the allowed memory.

% [2.19]
A measured run ends in one of a few outcomes.
It completes.
The OOM killer kills it because the cap is exceeded.
It aborts because the buffer pool is exhausted (Section~\ref{sec:bg-nes}).
It exceeds the time limit.
Or, rarely, it crashes on the restore path, an outcome confined to cells that already fail due to memory issues (Section~\ref{sec:rq3}).
The full mapping of each outcome to its process exit code is given later, in Table~\ref{tab:outcomes}.

% [2.20] The statistical tests are introduced in the evaluation itself
% (Section 5.1, [7.1.8]), next to the measurements they are applied to.

These terms complete the background.
The next chapter states the problem of this thesis and the questions the evaluation answers.
